% ============================================================
\chapter{Riemannian Manifolds}
\label{ch:manifolds}
% ============================================================

In 1854, Bernhard Riemann, only 27 years old, delivered his habilitation
lecture at G\"ottingen before a jury that included Gauss himself. The title:
\emph{\"Uber die Hypothesen, welche der Geometrie zu Grunde liegen} ---
``On the hypotheses which lie at the foundations of geometry.'' In one hour,
Riemann overturned the foundations of geometry: he proposed replacing rigid
Euclidean space with \emph{curved spaces}, where the notion of distance
varies from point to point. Gauss, it is said, was profoundly moved.

Riemann's central idea is elegant: instead of defining a global distance,
one equips each point with an \emph{inner product} on its tangent space.
From this infinitesimal data, one can reconstruct distances, angles, areas,
and even the curvature of space. It is this structure --- a manifold endowed
with such a field of inner products --- that we call a \emph{Riemannian
manifold}.

% ------------------------------------------------------------
\section{Riemannian metrics}
% ------------------------------------------------------------

Let us formalize Riemann's intuition. On a smooth manifold $M$, each point
$p$ possesses a tangent space $T_pM$, which is a vector space of dimension
$n$. A Riemannian metric consists of choosing an inner product on each of
these tangent spaces, in a smooth way.

\begin{definition}[Riemannian metric]
Let $M$ be a smooth manifold of dimension $n$. A \emph{Riemannian metric} on $M$
is the assignment, for each point $p \in M$, of an inner product
$g_p : T_pM \times T_pM \to \R$ varying smoothly: for all smooth vector fields
$X, Y \in \mathfrak{X}(M)$, the function $p \mapsto g_p(X_p, Y_p)$ is $C^\infty$.

The pair $(M, g)$ is called a \emph{Riemannian manifold}.
\end{definition}

In local coordinates $(x^1, \ldots, x^n)$, the metric is written:
\[
    g = g_{ij} \, dx^i \otimes dx^j, \qquad
    g_{ij} = g\!\left(\frac{\partial}{\partial x^i}, \frac{\partial}{\partial x^j}\right).
\]
The matrix $(g_{ij}(p))$ is symmetric positive definite at every point $p$.

\begin{theorem}[Existence of Riemannian metrics]
Every paracompact smooth manifold admits a Riemannian metric.
\end{theorem}

\begin{proof}
Let $\{(U_\alpha, \varphi_\alpha)\}$ be an atlas of $M$ and $\{\rho_\alpha\}$ a
partition of unity subordinate to $\{U_\alpha\}$. On each chart domain $U_\alpha$,
define the pullback Euclidean metric
$g_\alpha = \varphi_\alpha^*(\delta_{ij}\, dx^i \otimes dx^j)$.
Then $g = \sum_\alpha \rho_\alpha \, g_\alpha$ is a Riemannian metric on $M$,
since a convex combination of inner products is again an inner product.
\end{proof}

\begin{remark}
Uniqueness fails: the space of Riemannian metrics on $M$ is an infinite-dimensional
open convex cone in $\Gamma(S^2 T^*M)$.
\end{remark}

\section{Fundamental examples}

\begin{example}[Euclidean space]
The space $\R^n$ with the standard metric $g = \delta_{ij}\, dx^i \otimes dx^j$
is the simplest Riemannian manifold. We have $g_{ij} = \delta_{ij}$.
\end{example}

\begin{example}[The sphere $S^n$]
\label{ex:sphere}
The unit sphere $S^n = \{x \in \R^{n+1} : \norm{x} = 1\}$ carries the induced
metric from the inclusion $\iota : S^n \hookrightarrow \R^{n+1}$:
$g_{S^n} = \iota^* g_{\R^{n+1}}$.

In spherical coordinates on $S^2$, with $(\theta, \phi) \in (0,\pi) \times (0,2\pi)$:
\[
    g_{S^2} = d\theta^2 + \sin^2\!\theta \, d\phi^2.
\]
The metric matrix is
$\displaystyle (g_{ij}) = \begin{pmatrix} 1 & 0 \\ 0 & \sin^2\!\theta \end{pmatrix}$,
with determinant $\det(g_{ij}) = \sin^2\!\theta$.
\end{example}

\begin{example}[Hyperbolic space $\mathbb{H}^n$]
\label{ex:hyperbolic}
The \emph{upper half-space model}: $\mathbb{H}^n = \{(x^1, \ldots, x^n)
\in \R^n : x^n > 0\}$ with metric:
\[
    g_{\mathbb{H}^n} = \frac{(dx^1)^2 + \cdots + (dx^n)^2}{(x^n)^2}.
\]
For $n = 2$: $g = \frac{dx^2 + dy^2}{y^2}$ with $(x, y) \in \R \times \R_{>0}$.

The \emph{Poincar\'e disk model}: $\mathbb{D}^n = \{x \in \R^n : \norm{x} < 1\}$
with metric:
\[
    g_{\mathbb{D}^n} = \frac{4\left((dx^1)^2 + \cdots + (dx^n)^2\right)}{(1 - \norm{x}^2)^2}.
\]
These two models are isometric.
\end{example}

\begin{example}[Flat torus]
\label{ex:torus}
The torus $\mathbb{T}^n = \R^n / \Z^n$ inherits the standard Euclidean metric
by passing to the quotient. It is a compact flat manifold (zero curvature).
More generally, for a lattice $\Lambda \subset \R^n$, the torus $\R^n/\Lambda$
is a flat Riemannian manifold.
\end{example}

\begin{example}[Lie groups and bi-invariant metrics]
\label{ex:lie}
Let $G$ be a Lie group and $\mathfrak{g} = T_eG$ its Lie algebra. An inner
product $\inner{\cdot}{\cdot}$ on $\mathfrak{g}$ defines a left-invariant metric by:
\[
    g_a(u, v) = \inner{(dL_{a^{-1}})_a\, u}{(dL_{a^{-1}})_a\, v},
    \quad u, v \in T_aG,
\]
where $L_a : G \to G$ is left translation.

The metric is \emph{bi-invariant} if it is also right-invariant, which is equivalent to:
\[
    \inner{\operatorname{ad}_X Y}{Z} + \inner{Y}{\operatorname{ad}_X Z} = 0,
    \quad \forall\, X, Y, Z \in \mathfrak{g}.
\]
Every compact Lie group admits such a metric (for compact semisimple groups,
the negative of the Killing form works).
\end{example}

\section{Length, distance, and volume}

\begin{definition}[Length of a curve]
Let $\gamma : [a,b] \to M$ be a piecewise smooth curve. Its \emph{length} is:
\[
    L(\gamma) = \int_a^b \norm{\dot{\gamma}(t)} \, dt
    = \int_a^b \sqrt{g_{\gamma(t)}(\dot{\gamma}(t), \dot{\gamma}(t))} \, dt.
\]
\end{definition}

\begin{definition}[Riemannian distance]
The \emph{Riemannian distance} between two points $p, q \in M$ is:
\[
    d(p,q) = \inf \{ L(\gamma) : \gamma \text{ piecewise smooth curve from } p \text{ to } q \}.
\]
\end{definition}

\begin{theorem}[Metric space structure]
The distance $d$ defines a metric space structure on $M$ whose topology coincides
with the manifold topology.
\end{theorem}

\begin{proof}
The separation property $d(p,q) > 0$ for $p \neq q$ follows from working in
local coordinates, where the Riemannian metric is equivalent to the Euclidean
metric, so any curve joining distinct points has strictly positive length.
The triangle inequality is immediate from concatenation of curves. Topological
compatibility is proved by comparing metric balls with coordinate balls.
\end{proof}

\begin{definition}[Volume element]
The \emph{Riemannian volume form} is the $n$-form defined in oriented local
coordinates by:
\[
    dV_g = \sqrt{\det(g_{ij})} \, dx^1 \wedge \cdots \wedge dx^n.
\]
The volume of a domain $\Omega \subset M$ is $\operatorname{Vol}(\Omega) = \int_\Omega dV_g$.
\end{definition}

\begin{keyformulas}[title=\textbf{Volume formulas for model spaces}]
\begin{align*}
    \operatorname{Vol}(S^n(r)) &= \frac{2\pi^{(n+1)/2}}{\Gamma\!\left(\frac{n+1}{2}\right)} r^n, \\
    \operatorname{Vol}(B_{\R^n}(0,r)) &= \frac{\pi^{n/2}}{\Gamma\!\left(\frac{n}{2}+1\right)} r^n, \\
    \operatorname{Vol}(B_{\mathbb{H}^n}(p,r)) &= \omega_{n-1} \int_0^r \sinh^{n-1}(t)\, dt.
\end{align*}
\end{keyformulas}

\section{Isometries and conformal maps}

\begin{definition}[Isometry]
Let $(M, g)$ and $(N, h)$ be Riemannian manifolds. A diffeomorphism
$\varphi : M \to N$ is an \emph{isometry} if $\varphi^* h = g$, i.e.:
\[
    h_{\varphi(p)}(d\varphi_p(u), d\varphi_p(v)) = g_p(u, v),
    \quad \forall\, p \in M,\; u, v \in T_pM.
\]
\end{definition}

\begin{theorem}[Isometry group -- Myers--Steenrod]
The isometry group $\operatorname{Isom}(M, g)$ of a Riemannian manifold is a
Lie group of dimension at most $\frac{n(n+1)}{2}$.
\end{theorem}

\begin{remark}
Equality $\dim \operatorname{Isom}(M,g) = \frac{n(n+1)}{2}$ holds if and only if
$(M, g)$ has constant sectional curvature.
\end{remark}

\begin{example}[Isometries of $S^n$]
$\operatorname{Isom}(S^n) = O(n+1)$, of dimension $\frac{n(n+1)}{2}$.
\end{example}

\begin{example}[Isometries of $\mathbb{H}^n$]
In the upper half-space model, for $n=2$, the orientation-preserving isometries
are the M\"obius transformations $z \mapsto \frac{az+b}{cz+d}$ with
$a, b, c, d \in \R$, $ad - bc = 1$, giving $\operatorname{Isom}^+(\mathbb{H}^2)
\cong PSL(2, \R)$.
\end{example}

\begin{definition}[Conformal map]
A diffeomorphism $\varphi : (M,g) \to (N,h)$ is \emph{conformal} if there exists
$\lambda \in C^\infty(M)$, $\lambda > 0$, such that $\varphi^* h = \lambda^2 g$.
\end{definition}

\section{Products and Riemannian coverings}

\begin{definition}[Product metric]
Let $(M_1, g_1)$ and $(M_2, g_2)$ be Riemannian manifolds. The \emph{Riemannian
product} is $(M_1 \times M_2, g_1 \oplus g_2)$, where:
\[
    (g_1 \oplus g_2)_{(p_1, p_2)}((u_1, u_2), (v_1, v_2))
    = (g_1)_{p_1}(u_1, v_1) + (g_2)_{p_2}(u_2, v_2).
\]
\end{definition}

\begin{definition}[Warped product]
Let $(B, g_B)$ be a base and $(F, g_F)$ a fiber. For $f \in C^\infty(B)$, $f > 0$,
the \emph{warped product} is $B \times_f F$ with metric:
\[
    g = g_B + f^2 \, g_F.
\]
\end{definition}

\begin{example}
The space $\R^n \setminus \{0\}$ in polar coordinates is a warped product
$\R_{>0} \times_r S^{n-1}$ with $g = dr^2 + r^2 g_{S^{n-1}}$.
\end{example}

\begin{proposition}[Riemannian covering]
If $\pi : \widetilde{M} \to M$ is a covering map and $(M, g)$ is a Riemannian
manifold, then there exists a unique metric $\tilde{g}$ on $\widetilde{M}$
such that $\pi$ is a local isometry: $\pi^* g = \tilde{g}$.
\end{proposition}

\section{Completeness}

\begin{definition}[Metric completeness]
$(M, g)$ is \emph{complete} if the metric space $(M, d)$ is Cauchy-complete,
i.e., every Cauchy sequence converges.
\end{definition}

\begin{theorem}[Hopf--Rinow]
\label{thm:hopf_rinow}
Let $(M, g)$ be a connected Riemannian manifold. The following are equivalent:
\begin{enumerate}
    \item $(M, d)$ is complete as a metric space.
    \item Closed bounded subsets are compact.
    \item There exists $p \in M$ such that $\exp_p$ is defined on all of $T_pM$.
    \item For every $p \in M$, $\exp_p$ is defined on all of $T_pM$.
\end{enumerate}
Moreover, if any of these conditions holds, then any two points of $M$ are
joined by a minimizing geodesic.
\end{theorem}

\begin{warning}[title=\textbf{Completeness vs.\ compactness}]
Completeness does \emph{not} imply compactness: $\R^n$ and $\mathbb{H}^n$ are
complete but noncompact. Conversely, every compact Riemannian manifold is complete.
\end{warning}

\section{Coordinate computations}

\begin{example}[Polar coordinates on $\R^2$]
The coordinate change $x = r\cos\theta$, $y = r\sin\theta$ gives:
\[
    g = dx^2 + dy^2 = dr^2 + r^2\, d\theta^2.
\]
Hence $(g_{ij}) = \begin{pmatrix} 1 & 0 \\ 0 & r^2 \end{pmatrix}$ and
$\sqrt{\det g} = r$.
\end{example}

\begin{example}[Fubini--Study metric on $\C P^1$]
Identifying $\C P^1 \cong S^2$, in inhomogeneous coordinates $z = x + iy$:
\[
    g_{FS} = \frac{4\,(dx^2 + dy^2)}{(1 + x^2 + y^2)^2}
    = \frac{4\, |dz|^2}{(1 + |z|^2)^2}.
\]
This is the constant curvature $K = 1$ metric on $S^2$.
\end{example}

\section{Functorial operations on metrics}

\begin{proposition}[Pullback metric]
If $f : M \to N$ is an immersion and $(N, h)$ is Riemannian, then
$g = f^* h$ is a Riemannian metric on $M$ (the \emph{induced metric}).
\end{proposition}

\begin{proposition}[Musical isomorphisms]
The metric induces canonical isomorphisms:
\begin{align*}
    \flat : TM &\to T^*M, & v &\mapsto g(v, \cdot), \\
    \sharp : T^*M &\to TM, & \omega &\mapsto \text{the unique } v \text{ with }
    g(v, \cdot) = \omega.
\end{align*}
In coordinates: $(v^\flat)_i = g_{ij} v^j$ and $(\omega^\sharp)^i = g^{ij} \omega_j$,
where $(g^{ij})$ is the inverse matrix of $(g_{ij})$.
\end{proposition}

\section{Exercises}

\begin{exercise}
Show that the hyperbolic metric $g = \frac{dx^2 + dy^2}{y^2}$ on $\mathbb{H}^2$
gives the distance formula:
\[
    d((x_1, y_1), (x_2, y_2))
    = \operatorname{arcosh}\!\left(1 + \frac{(x_1-x_2)^2 + (y_1-y_2)^2}{2 y_1 y_2}\right).
\]
\end{exercise}

\begin{exercise}
Compute the volume of the sphere $S^2(r)$ of radius $r$ by integrating the volume
form in spherical coordinates. Verify that $\operatorname{Vol}(S^2(r)) = 4\pi r^2$.
\end{exercise}

\begin{exercise}
Let $G = SO(3)$ with the bi-invariant metric $g(X, Y) = -\frac{1}{2}\tr(XY)$
for $X, Y \in \mathfrak{so}(3)$. Show that this metric makes $SO(3)$ a space
of constant sectional curvature $K = 1/4$.
\end{exercise}

\begin{exercise}
Show that the warped product $M = \R \times_{\cosh t} \R$ with
$g = dt^2 + \cosh^2(t)\, ds^2$ is isometric to an open subset of $\mathbb{H}^2$.
\end{exercise}

\begin{exercise}
Let $(M, g)$ be a Riemannian manifold and $\lambda \in C^\infty(M)$, $\lambda > 0$.
Set $\tilde{g} = \lambda^2 g$ (conformal change). Compute the relation between
the volume forms $dV_{\tilde{g}}$ and $dV_g$.
\end{exercise}
